% from ega4-20-fr.tex
\oldpage[IV-4]{251}
appelle le pseudo-morphisme de $X$ dans $Z$, relativement à $S$, \emph{composé} de $f$ et de $\omega$ et que l'on note $f\circ\omega$.
Si $g:Z\to T$ est un morphisme, il est immédiat que $g\circ(f\circ\omega)=(g\circ f)\circ\omega$.
\end{env}

\begin{env}[20.5.6]
\label{IV.20.5.6}
Supposons $Y$ \emph{séparé} sur $S$, et soit $\omega$ un pseudo-morphisme de $X$ dans $Y$ relativement à $S$.
Si $f:X'\to X$ est un $S$-morphisme tel que $f^{-1}(\operatorname{dom}_S(\omega))$ soit \emph{universellement schématiquement dense dans $X'$ relativement à $S$}, il résulte de \sref{IV.20.3.2} que le pseudo-$S$-morphisme $\omega\circ f$ a un domaine universellement schématiquement dense relativement à $S$, donc \sref{IV.20.5.3} peut être considéré comme un pseudo-morphisme \emph{relativement à $S$}.
Lorsque $X'$ est plat et localement de présentation finie sur $S$, la condition que $f^{-1}(\operatorname{dom}_S(\omega))$ soit universellement schématiquement dense relativement à $S$ équivaut encore à dire que pour tout $s\in S$, $(f^{-1}(\operatorname{dom}_S(\omega)))_s$ (notation de \sref{IV.11.10.10}) est schématiquement dense dans $X'_s$, ou encore, en notant $f_s:X'_s\to X_s$ le morphisme déduit de $f$ par changement de base, que l'image réciproque par $f_s$ de $(\operatorname{dom}_S(\omega))_s$ soit schématiquement dense dans $X'_s$.
Cette dernière condition sera en particulier remplie si, pour tout $s\in S$, $X_s$, $X'_s$ et $f_s$ vérifient l'une des trois conditions~(i),~(ii),~(iii) de \sref{IV.20.3.5}.
\end{env}

\begin{env}[20.5.7]
\label{IV.20.5.7}
Supposons maintenant que $X$ et $X'$ soient tous deux des $S$-préschémas \emph{plats et localement de présentation finie sur $S$}, et que $f:X'\to X$ soit un $S$-morphisme \emph{plat} (ou, ce qui revient au même \sref{IV.11.3.10}, que pour tout $s\in S$, $f_s:X'_s\to X_s$ soit un morphisme plat).
Alors, pour tout ouvert $U$ de $X$ et tout ouvert $V\subset U$ universellement schématiquement dense dans $U$ relativement à $S$, il résulte de \sref{IV.11.10.5} et \sref{IV.11.10.9} que $f^{-1}(V)$ est universellement schématiquement dense dans $f^{-1}(U)$ relativement à $S$.
Pour tout pseudo-morphisme $\omega$ de $X$ dans $Y$ relativement à $S$, il résulte de \sref{IV.20.3.4} que le pseudo-$S$-morphisme $\omega\circ f$ est défini et est un pseudo-morphisme de $X'$ dans $Y$ relativement à $S$, \emph{même lorsqu'on ne suppose pas $Y$ séparé sur $S$}.
On en déduit que dans ce cas, pour tout $S$-morphisme $g:Y\to Z$, $(g\circ\omega)\circ f$ est encore défini et égal à $g\circ(\omega\circ f)$ (avec les définitions de \sref{IV.20.5.1}), et est donc un pseudo-morphisme \emph{relativement à $S$}.
\end{env}

\begin{env}[20.5.8]
\label{IV.20.5.8}
Soient $X$ un $S$-préschéma, $S'\to S$ un morphisme quelconque, $X'=X_{(S')}$, $g:X'\to X$ la projection canonique.
Alors, pour tout ouvert $U$ de $X$ et tout ouvert $V\subset U$ universellement schématiquement dense dans $U$ relativement à $S$, $V'=g^{-1}(V)$ est universellement schématiquement dense dans $U'=g^{-1}(U)$ relativement à $S'$ par définition \sref{IV.11.10.8}.
Soit alors $\omega$ un pseudo-morphisme de $X$ dans un $S$-préschéma $Y$ relativement à $S$ ; si $u_1$, $u_2$ sont des $S$-morphismes de la classe $\omega$, définis respectivement dans des ouverts $U_1$, $U_2$ de $X$, universellement schématiquement denses dans $X$ relativement à $S$, il résulte de ce qui précède que les morphismes $u'_1=u_1\circ(g|g^{-1}(U_1))$ et $u'_2=u_2\circ(g|g^{-1}(U_2))$ coïncident dans un ouvert $U'_3$ universellement schématiquement dense relativement à $S'$.
Or, si $Y'=Y_{(S')}$ et si $h:Y'\to Y$ est la projection canonique, $u'_1$ et $u'_2$ se factorisent canoniquement en $u'_1=h\circ v_1$, $u'_2=h\circ v_2$, et $v_1$ et $v_2$ sont deux $S'$-morphismes dans $Y'$ qui coïncident dans $U'_3$.
On voit donc que lorsque $u_1$ parcourt la classe $\omega$, les $S'$-morphismes $v_1$ correspondants appartiennent à un même pseudo-morphisme relativement à $S'$, dit \emph{image réciproque} de $\omega$ par le morphisme changement de base $S'\to S$ et noté $\omega_{(S')}$.
Il est clair que si $S''\to S'$ est un second morphisme, on a $(\omega_{(S')})_{(S'')}=\omega_{(S'')}$ (pour le morphisme changement de base composé $S''\to S'\to S$).
\end{env}

